\title{FlashAttention-3: Fast and Accurate Attention with Asynchrony and Low-precision}

\begin{abstract}
Attention mechanisms constitute the fundamental computational layer within Transformer architectures, acting as the primary performance constraint for scaling large language models and enabling extensive context processing. \textsc{FlashAttention} previously introduced techniques for accelerating attention calculations on GPUs by reducing memory access overhead. Nonetheless, existing implementations have yet to fully harness the advanced functionalities offered by emerging GPU hardware; for instance, \textsc{FlashAttention-2} utilizes only 35\% of the computational capacity on the H100 GPU. In this work, we present three strategies to enhance attention computation specifically on Hopper GPUs: we (1) capitalize on the asynchrony between Tensor Cores and TMA to overlap data transfer and computation via warp specialization, (2) interleave block-level matrix multiplications with softmax evaluations, and (3) apply block quantization and incoherent processing, leveraging hardware acceleration for FP8 precision. Employing these methods, our approach, termed \textsc{FlashAttention-3}, delivers 1.5-2.0$\times$ acceleration on H100 GPUs, with FP16 performance peaking at 740 TFLOPs/s (representing 75\% utilization), and FP8 execution nearing 1.2 PFLOPs/s. Our evaluation further establishes that FP8 \textsc{FlashAttention-3} attains 2.6$\times$ lower numerical error compared to conventional FP8 attention baselines.
\end{abstract}

\section{Introduction}
\label{sec:intro}

Within the Transformer architecture~\citep{vaswani2017attention}, the core computational challenge lies in the attention mechanism, as calculating self-attention between queries and keys exhibits quadratic complexity with respect to sequence length. Extending attention to accommodate longer sequences offers the potential for advancements, such as enhanced modeling and reasoning across extensive documents~\citep{guo2021longt5,shaham2022scrolls,peng2023yarn} and handling multiple files within large code repositories~\citep{roziere2023code, li2023starcoder}, as well as supporting additional modalities like high-resolution images~\citep{chen2022scaling}, audio~\citep{gulati2020conformer}, and video~\citep{ho2022video}. Further applications include tracking user interactions with extended histories~\citep{sun2019bert4rec} and orchestrating agent workflows over prolonged time frames~\citep{yao2022react}. These prospects have motivated considerable research into accelerating attention for long-context scenarios, either through approximation techniques~\citep{katharopoulos2020transformers,choromanski2020rethinking,
tay2020efficient}, software engineering approaches (\citep{rabe2021self,dao2022flashattention,kwon2023efficient}), or by adopting alternative model architectures~\citep{peng2023rwkv,sun2023retentive,gu2023mamba}.

This study extends the foundational contributions by \citet{dao2022flashattention}, who designed exact-attention algorithms by explicitly considering the GPU's execution paradigms and low-level hardware features in their algorithmic architecture. The \textsc{FlashAttention} approach presented in \citep{dao2022flashattention} introduced an innovative tiling mechanism for accelerating attention by fusing all computational stages of attention into a unified GPU kernel, thus obviating the need for frequent access to slower global memory between stages. Building upon this, \citet{dao2023flashattention2} proposed \textsc{FlashAttention-2}, in which the algorithm is further adapted to parallelize across the sequence length, with the innermost computations in the forward pass distributed over block partitions of the key and value matrices, leading to enhanced load balancing and utilization on the GPU. Nevertheless, we note that despite these improvements, \textsc{FlashAttention-2} does not fully exploit the computational power of recent GPUs, performing notably worse than state-of-the-art matrix multiplication (GEMM) implementations—realizing only 35\% utilization as opposed to 80–90\% on the Hopper H100 GPU. This discrepancy can be partly explained by differences at the implementation level, such as the reliance on Ampere-optimized instructions over those specific to the Hopper architecture for Tensor Core operations. Contemporary efforts, including ThunkerKitten~\citep{spector2024thunder} and cuDNN 9~\citep{cudnn9}, demonstrate that the adoption of Hopper-tailored instructions combined with tile-based abstractions can significantly enhance the efficiency of attention mechanisms and reduce implementation complexity.

At a more basic level, the algorithm used by \textsc{FlashAttention-2} operates within a streamlined synchronous framework and does not actively incorporate asynchronous operations or low-precision data formats in its architecture. Asynchronous behavior emerges due to hardware optimizations, which prioritize high-impact computations in machine learning workflows: specialized units handle matrix multiplication (Tensor Cores) and memory transfers (Tensor Memory Accelerator -- TMA), functioning separately from conventional CUDA cores that execute control logic and perform both integer and floating point calculations. The use of reduced numerical precision, such as FP8 available in Hopper and FP4 in Blackwell, follows the established practice observed with FP16 (as introduced with Pascal in 2017) and BF16 (featured in Ampere in 2020), effectively increasing computational throughput two- or fourfold within identical power and silicon footprint constraints. The features specific to Hopper regarding these developments are discussed in detail in \cref{subsec:hardware}. The primary engineering hurdle lies in modifying the \textsc{FlashAttention-2} framework to leverage these hardware innovations: asynchronous execution necessitates the concurrent calculation of matmul and softmax operations, despite the dependency chain between their outputs, and employing lower precision mandates careful strategies to prevent excessive quantization artifacts, which are particularly problematic when dealing with rare features in LLMs~\citep{dettmers2208llm, sun2024massive}.

To address these objectives, we introduce \textsc{FlashAttention-3}, which integrates three novel approaches aimed at enhancing efficiency on modern GPU platforms:\footnote{Our experiments are presented using NVIDIA's Hopper architecture as the reference. Nevertheless, the algorithm remains applicable to any GPU that supports advanced asynchronous operations and low-precision computation features.}

\begin{enumerate}

\item \textbf{Producer-Consumer asynchrony:} We introduce a warp-specialized software pipelining approach that leverages asynchronous data movement and Tensor Core execution. By assigning the roles of data producers and consumers to distinct warps, this mechanism enhances the algorithm's latency tolerance for memory accesses and instruction dispatch.
\item \textbf{Concealing softmax latency within asynchronous block-wise GEMMs:} We orchestrate the execution of low-throughput non-GEMM operations in softmax (e.g., floating-point multiply-add and exponentiation) to overlap with asynchronous WGMMA GEMM instructions. To achieve this, we refactor the \textsc{FlashAttention-2} algorithm, alleviating certain sequential dependencies between softmax computations and GEMM steps. For instance, in our two-stage variant, while softmax operates on one block of the score matrix, the asynchronous proxy executes WGMMA for the subsequent block.
\item \textbf{Hardware-accelerated GEMM at reduced precision:} We modify the forward algorithm to utilize FP8 Tensor Cores in GEMM computations, resulting in a substantial increase—nearly twofold—in observed TFLOPs/s. This adaptation requires resolving the layout differences for WGMMA with respect to the arrangement of FP32 accumulator and FP8 operand matrices in memory. To compensate for the reduced numerical accuracy at FP8, we employ block quantization and incoherent computation strategies.
\end{enumerate}

For empirical validation, we conducted benchmarks of \textsc{FlashAttention-3} on an H100 SXM5 GPU across a range of settings. Our results indicate that (1) FP16 offers a performance increase of 1.5-2.0$\times$ over \textsc{FlashAttention-2} in the forward computation—attaining speeds up to 740 TFLOPs/s—and achieves 1.5-1.75$\times$ improvement in the backward computation, (2) FP8 approaches 1.2 PFLOPs/s throughput, and (3) for extended sequence lengths, FP16 exhibits superior performance and FP8 shows competitive results\footnote{Specifically, \textsc{FlashAttention-3} FP8 is advantageous with head dimension 64; for head dimensions 128 and 256 it performs comparably to state-of-the-art without causal masking, and less favorably when causal masking is employed.} relative to NVIDIA's cuDNN library's attention kernel. Additionally, we confirm that the FP16 version of \textsc{FlashAttention-3} maintains numerical errors on par with those in \textsc{FlashAttention-2}, and performs better than the conventional attention algorithm, as it retains intermediate calculations (such as softmax scaling) in FP32 precision. Furthermore, FP8 \textsc{FlashAttention-3}, utilizing block quantization and incoherent execution, achieves 2.6$\times$ higher accuracy than traditional attention employing per-tensor quantization when faced with features exhibiting significant outliers.

\textsc{FlashAttention-3} is released under an open-source, permissive license\footnote{\textsc{FlashAttention-3}
  is available at \texttt{https://github.com/Dao-AILab/flash-attention}}. Additionally, we intend to incorporate the implementation into both PyTorch and Hugging Face libraries, aiming to maximize its accessibility for the research and development community.

\section{Background: Multi-Head Attention and GPU Characteristics}
\label{sec:background}

\subsection{Multi-Head Attention}
\label{subsec:multi_head_attn}

Let $\mathbf{Q}, \mathbf{K}, \mathbf{V} \in \mathbb{R}^{N \times d}$ be the query, key and value input sequences associated to a single head, where $N$ is the sequence length and $d$ is the head dimension. Then the attention output $\mathbf{O}$ is computed as:

\begin{equation*}
  \mathbf{S} = \alpha \mathbf{Q} \mathbf{K}^\top \in \mathbb{R}^{N \times N}, \quad \mathbf{P} = \mathrm{softmax}(\mathbf{S}) \in \mathbb{R}^{N \times N}, \quad \mathbf{O} = \mathbf{P}\mathbf{V} \in \mathbb{R}^{N \times d},
\end{equation*}

Here, the $\mathrm{softmax}$ function operates across the rows, and commonly, $\alpha$ is chosen to be $1/\sqrt{d}$ to act as the scale parameter. To address potential numerical issues arising from exponentiation, $\mathrm{rowmax}(\mathbf{S})$ is subtracted from $\mathbf{S}$ prior to applying the exponential. In the case of multi-head attention (MHA), individual heads maintain separate query, key, and value projections. This process is efficiently parallelized over numerous heads and batches, culminating in the generation of the complete output tensor.

Now let $\phi$ be a scalar loss function and let $\mathbf{d}(-) = \partial \phi / \partial (-)$ be notation for the gradient. Given the output gradient $\mathbf{dO} \in \mathbb{R}^{N \times d}$, we compute $\mathbf{dQ}$, $\mathbf{dK}$, and $\mathbf{dV}$ according to the chain rule as follows:

\begin{align*}
  \mathbf{dV} &= \mathbf{P}^\top \mathbf{dO} \in \mathbb{R}^{N \times d} \\
  \mathbf{dP} &= \mathbf{dO} \mathbf{V}^\top \in \mathbb{R}^{N \times N} \\
  \mathbf{dS} &= \mathrm{dsoftmax} (\mathbf{dP}) \in \mathbb{R}^{N \times N} \\
  \mathbf{dQ} &= \alpha \mathbf{dS} \mathbf{K} \in \mathbb{R}^{N \times d} \\
  \mathbf{dK} &= \alpha \mathbf{dS}^\top \mathbf{Q} \in \mathbb{R}^{N \times d},
\end{align*}

In this context, $\mathbf{d}s = (\mathrm{diag}(p) - p p^\top)\mathbf{d}p$ holds when $p = \mathrm{softmax}(s)$, treating $p$ as dependent on the vector $s$. We denote the row-wise application of this relation as $\mathrm{dsoftmax}(\mathbf{dP})$. For the backward propagation in MHA, this calculation similarly lends itself to parallelization over both heads and batches.

\subsection{GPU hardware characteristics and execution model}
\label{subsec:hardware}

We outline the features of the GPU execution model pertinent to \textsc{FlashAttention-3}, emphasizing the NVIDIA Hopper architecture as a specific example within this context.

\paragraph{Memory hierarchy:} On GPUs, the memory system is structured as a series of hierarchical data locations, where greater capacity typically comes at the expense of reduced bandwidth (\cref{tab:gpu-hierarchy})\footnote{\citet{luo2024benchmarking} notes that shared memory offers a bandwidth of 128 bytes per clock cycle for every SM, which, when scaled by 132 SMs and a boost clock of 1830 MHz, yields the aggregate bandwidth.}. At the highest level, global memory (GMEM)—frequently referred to as HBM—serves as off-chip DRAM that is universally accessible by all streaming multiprocessors (SMs). Accesses to GMEM are automatically buffered in an on-chip L2 cache. Each SM, in turn, is equipped with a small, highly banked, on-chip memory that is explicitly managed by programmers, known as shared memory (SMEM). At the most granular level is the register file resident within each SM.

\paragraph{Thread hierarchy:} In the GPU programming paradigm, execution units are systematically grouped into logical entities known as threads. The structure of this hierarchy, ordered from the most granular to the broadest level, includes individual threads, warps (each consisting of 32 threads), warpgroups (made up of 4 consecutive warps), threadblocks (also termed cooperative thread arrays or CTAs), threadblock clusters (specific to Hopper architectures), and finally, grids.

The two hierarchies exhibit a strong interdependence. Threads belonging to a given CTA are assigned to execute together on a single SM, while CTAs grouped within the same cluster are concurrently executed on one GPC. All threads in a CTA have direct access to SMEM, but each individual thread possesses up to 256 dedicated registers (RMEM), which are not shared.

\paragraph{Asynchrony and warp-specialization:}

GPUs function as throughput-oriented processors, utilizing concurrent and asynchronous operations to mask the delays associated with memory access and execution. In Hopper, the Tensor Memory Accelerator (TMA) serves as specialized hardware for asynchronous memory transfers between GMEM and SMEM \cite[\S7.29]{cuda}. Additionally, Hopper’s architecture distinguishes itself from predecessors like Ampere by implementing asynchronous capabilities within the Tensor Core, which is accessible through the warpgroup-wide WGMMA instruction \cite[\S9.7.14]{ptx}; this allows the Tensor Core to fetch inputs directly from shared memory.

Enabling asynchrony at the hardware level makes it possible to implement warp-specialized kernels, in which the warps within a CTA are assigned solely to either data movement (producer) or computation (consumer) tasks. This separation, in a general sense, enhances the compiler's capacity to construct highly efficient instruction schedules \citep{warp-specialization-2011}. Further, Hopper introduces the ability to dynamically redistribute registers among different warpgroups through the \verb|setmaxnreg| mechanism \cite[\S9.7.17.1]{ptx}. As a result, warps engaged in MMA operations can be allocated a larger portion of RMEM compared to those restricted to issuing TMA commands, which only require a single thread.

\paragraph{Low-precision number formats:}
\label{sec:low-precision-gpu}
Contemporary GPUs incorporate dedicated hardware components to facilitate efficient low-precision arithmetic. As an illustration, the WGMMA instruction leverages the FP8 Tensor Cores available on the Hopper architecture, achieving double the throughput per SM relative to computations utilizing FP16 or BF16.

Effectively using FP8 WGMMA requires an awareness of the specific layout restrictions associated with its inputs. Consider a GEMM operation computing $A \times B^{\top}$, where $A$ is an $M\times K$ matrix and $B$ is $N\times K$. We describe an operand as \emph{mn-major} when it is stored contiguously along the $M$ or $N$ (outer) dimension, and as \emph{k-major} when contiguity occurs in the $K$ (inner) dimension. For FP16 WGMMA, both mn-major and k-major formats are permissible for SMEM-resident operands. Conversely, FP8 WGMMA imposes stricter requirements, accepting solely k-major inputs. This constraint becomes particularly challenging in applications such as attention mechanisms, where the intent is to merge consecutive GEMM operations into one kernel: mismatched layouts between FP32 accumulators and FP8 operands complicate direct usage of FP8 WGMMAs when dependencies exist between them.

When considering attention mechanisms, these layout constraints necessitate specific adjustments in the development of an FP8 algorithm, detailed in \cref{sec:algofp8}.

\subsection{Standard Attention and Flash Attention}

As described in \citet{dao2022flashattention}, \textbf{standard attention} on the GPU is typically implemented by explicitly storing the intermediate matrices $\mathbf{S}$ and $\mathbf{P}$ in high-bandwidth memory (HBM). The key contribution of \textsc{FlashAttention} was to exploit a localized variant of softmax reduction, effectively eliminating the need for costly memory accesses to these intermediates and enabling the attention computation to be combined into a single GPU kernel. This local softmax approach appears in lines \ref{code-ws:softmax_start}-\ref{code-ws:softmax_end} of the consumer mainloop within \cref{alg:flash3_wgmma_ws_only}, in concert with the rescaling operations performed on blocks of $\mathbf{O}$. A straightforward derivation that shows this method accurately yields $\mathbf{O}$ is available in \cite[\S 2.3.1]{dao2023flashattention2}.

\section{FlashAttention-3: Algorithm}
\label{sec:algo}

This section details the \textsc{FlashAttention-3} algorithm, with an emphasis on the forward computation; the procedure for the backward pass is outlined in~\cref{sec:algo_ws_bwd}. Initially, we illustrate the integration of warp-specialization combined with a circular SMEM buffer into the foundational \textsc{FlashAttention-2} algorithm. Subsequently, we discuss how leveraging the asynchronous capabilities of WGMMA enables us to establish a two-stage pipeline where GEMM and softmax operations are overlapped. In the final part, we outline the necessary adaptations for FP8, addressing both layout compatibility and precision enhancements through block quantization and incoherent processing.

\subsection{Producer-Consumer asynchrony through warp-specialization and pingpong scheduling}
\label{sec:algo_ws}

\paragraph{Warp-specialization}
Similar to \textsc{FlashAttention-2}, the forward computation in \textsc{FlashAttention-3} lends itself to significant parallelism across the batch dimension, the number of attention heads, and the length of the query sequence. Therefore, it is sufficient to provide a CTA-level overview in which each algorithmic unit processes a tile $\mathbf{Q}_i$ from the query matrix and produces the respective output tile $\mathbf{O}_i$. For clarity, we describe here the warp-specialization approach utilizing a circular SMEM buffer but omitting GEMM-softmax overlap. Let $d$ represent the dimension of each attention head, and $N$ denote the total sequence length. By selecting a query block size $B_r$, the matrix $\mathbf{Q}$ is partitioned into $T_r = \lceil \frac{N}{B_r} \rceil$ blocks: $\mathbf{Q}_1, \dots, \mathbf{Q}_{T_r}$.

\begin{algorithm}[H]
    \caption{\small\label{alg:flash3_wgmma_ws_only}\textsc{FlashAttention-3} forward pass \textbf{excluding} intra-consumer overlap -- CTA perspective}
    \begin{algorithmic}[1]
\REQUIRE Inputs: $\mathbf{Q}_i \in \mathbb{R}^{B_r \times d}$ and $\mathbf{K}, \mathbf{V} \in \mathbb{R}^{N \times d}$ located in HBM, a key block size $B_c$, and $T_c = \lceil \frac{N}{B_c} \rceil$ blocks.
\STATE Set up pipeline controller to handle barrier synchronization using an $s$-stage circular shared memory buffer.
\IF {in the producer warpgroup}
\STATE Free the predetermined number of registers for use.
\STATE Execute transfer of $\mathbf{Q}_i$ from HBM into shared memory.
\STATE Once transfer finishes, signal to consumer that $\mathbf{Q}_i$ is available.
\FOR{$j$ from $0$ to $T_c - 1$}
    \STATE Await that the $(j\,\%\,s)$ buffer stage has been processed by consumers.
    \STATE Trigger loading of $\mathbf{K}_j$ and $\mathbf{V}_j$ from HBM into the $(j\,\%\,s)$ buffer stage in shared memory.
    \STATE On finishing, confirm to consumers that both $\mathbf{K}_j$ and $\mathbf{V}_j$ are ready.
\ENDFOR
\ELSE \STATE Allocate registers based on the number of consumer warps.
\STATE Locally initialize $\mathbf{O}_i = (0) \in \mathbb{R}^{B_r \times d}$, $\ell_i = (0)$, and $m_i = (-\infty)$ with $\ell_i, m_i \in \mathbb{R}^{B_r}$.
\STATE Wait until $\mathbf{Q}_i$ appears in shared memory.
\FOR{$j$ from $0$ to $T_c - 1$}
\STATE Await completion of $\mathbf{K}_j$ load into shared memory.
\STATE Calculate $\mathbf{S}_i^{(j)} = \mathbf{Q}_i \mathbf{K}_j^T$ using SS-GEMM. Commit and synchronize. \label{alg:ws_only_gemm1}
\STATE Preserve previous $m_i$ as $m_i^{\mathrm{old}}$, then update: $m_i = \mathrm{max}(m_i^{\mathrm{old}}, \mathrm{rowmax}(\mathbf{S}_i^{(j)}))$. \label{code-ws:softmax_start}
\STATE Evaluate $\widetilde{\mathbf{P}}_i^{(j)} = \mathrm{exp}(\mathbf{S}_i^{(j)} - m_i)$ and refresh $\ell_i = \mathrm{exp}(m_i^{\mathrm{old}} - m_i) \ell_i + \mathrm{rowsum}(\widetilde{\mathbf{P}}_i^{(j)})$. \label{code-ws:softmax_end}
\STATE Hold until $\mathbf{V}_j$ finishes loading into shared memory.
\STATE Update $\mathbf{O}_i$ by $\mathbf{O}_i = \mathrm{diag}(\mathrm{exp}(m_i^{\mathrm{old}} - m_i))^{-1} \mathbf{O}_i + \widetilde{\mathbf{P}}_i^{(j)} \mathbf{V}_j$ using RS-GEMM. Commit and synchronize. \label{alg:ws_only_gemm2}
\STATE Mark the $(j\,\%\,s)$ buffer stage as released for producers.
\ENDFOR
\STATE Finalize output with $\mathbf{O}_i = \mathrm{diag}(\ell_i)^{-1} \mathbf{O}_i$ and calculate $L_i = m_i + \log(\ell_i)$.
\STATE Write out $\mathbf{O}_i$ and $L_i$ into HBM as the $i$th section of $\mathbf{O}$ and $L$.
\ENDIF
\end{algorithmic}
\end{algorithm}

In our deployment of \cref{alg:flash3_wgmma_ws_only} on Hopper, we employ \verb|setmaxnreg| to manage (de)allocation operations, utilize TMA to fetch data for $\mathbf{Q}_i$ as well as $\{ \mathbf{K}_j, \mathbf{V}_j \}_{0 \leq j < T_c}$, and rely on WGMMA for performing the GEMMs within the consumer main loop. The SS or RS prefix identifies whether the initial operand is obtained from shared memory or the register file, respectively. When analyzing the execution trajectory of \cref{alg:flash3_wgmma_ws_only}, it is important to recognize that TMA load instructions can be issued without blocking on prior loads, thanks to their asynchronous nature. Additionally, within the producer main loop, no synchronization waits are necessary during the initial $s$ iterations as the buffer undergoes its filling phase.

\paragraph{Pingpong scheduling} The asynchronous execution of WGMMA and TMA, combined with warp-specialization, enables the possibility of concurrently executing the softmax calculation for one warpgroup while another warpgroup performs GEMM. This parallelism is particularly advantageous when considering that, on contemporary hardware accelerators, operations other than matmul (such as exponentials) exhibit significantly less throughput than matmul computations. To illustrate, the H100 SXM5 GPU boasts a FP16 matmul capacity of 989 TFLOPS, but its throughput for special functions like exponential is only 3.9 TFLOPS\footnote{According to the CUDA programming guide, each streaming multiprocessor (SM) can carry out 16 special function operations per clock cycle. By multiplying 16 by the 132 SMs and the 1830 MHz clock frequency, one arrives at 3.9 TFLOPS for these functions.}—a stark contrast required for softmax. For the attention forward step in FP16 with a head dimension of 128, the volume of matmul FLOPS is 512 times higher than that required for exponentials, yet exponentials are processed at 256 times lower throughput, meaning they can occupy around 50\% of the total cycles taken by matmul. This imbalance is exacerbated with FP8, where matmul throughput is twice that of FP16, but the exponential operation rate remains unchanged.

Because the exponential operation is handled by a dedicated hardware unit, namely the multi-function unit, it is preferable to schedule this computation concurrently with the matmul tasks executed by the Tensor Cores. To achieve such overlap, we employ synchronization barriers (\texttt{bar.sync} instructions) that enforce an execution order: the GEMMs (specifically, GEMM1 -- $\mathbf{P} \mathbf{V}$ within the current iteration, and GEMM0 -- $\mathbf{Q} \mathbf{K}^\top$ for the upcoming iteration) assigned to warpgroup 1 are completed prior to initiating those for warpgroup 2. This allows warpgroup 1 to proceed with the softmax calculation while warpgroup 2 undertakes its GEMMs. Subsequently, their responsibilities alternate, so warpgroup 2 performs softmax while warpgroup 1 executes GEMMs; this strategy is referred to as ``pingpong'' scheduling. The scheduling approach is depicted in~\cref{fig:pingpong_scheduling}. While the real-world implementation of pingpong scheduling introduces some irregularities compared to the conceptual illustration, it nonetheless typically yields noticeable performance gains (for instance, increasing throughput from 570 TFLOPS to 620-640 TFLOPS in FP16 forward passes with a head dimension of 128 and sequence length of 8192).

\paragraph{Attention variants} In the cases of multi-query attention~\citep{shazeer2019fast} and grouped query attention~\citep{ainslie2023gqa}, our implementation mirrors that of \textsc{FlashAttention-2}, modifying tensor indexing so that $\mathbf{K}$ and $\mathbf{V}$ are not redundantly stored within HBM.

\subsection{Intra-warpgroup overlapping GEMMs and softmax}

It is possible to execute certain instructions from the softmax operation simultaneously with selected instructions from the GEMMs, even within a single warpgroup. Here, we outline a method to achieve this overlap.

In the attention algorithm, operations within the inner loop (main loop) have sequential dependencies that impede parallelization within a single iteration. For example, (local) softmax (lines \ref{code-ws:softmax_start} to \ref{code-ws:softmax_end}) relies on the output $\mathbf{S}_i^{(j)}$ of the first GEMM, while the second GEMM takes its result $\widetilde{\mathbf{P}}_i^{(j)}$ as an operand. Indeed, the wait statements in lines \ref{alg:ws_only_gemm1} and \ref{alg:ws_only_gemm2} of \cref{alg:flash3_wgmma_ws_only} serialize the execution of softmax and GEMMs. However, we can break these dependencies by pipelining across iterations through additional buffers in registers. Pursuing this idea, we propose the following two-stage\footnote{Note that the number of stages of the overlapping scheme is bounded by, but need not equal, the number $s$ of stages in the circular SMEM buffer.} GEMM-softmax pipelining algorithm:

\begin{algorithm}[H]
    \caption{\small\label{alg:flash3_wgmma}\textsc{FlashAttention-3} consumer warpgroup forward pass}
    \begin{algorithmic}[1]
      \REQUIRE  Matrices $\mathbf{Q}_i \in \mathbb{R}^{B_r \times d}$ and $\mathbf{K}, \mathbf{V} \in \mathbb{R}^{N \times d}$ reside in HBM, with a key block size $B_c$, yielding $T_c = \lceil \frac{N}{B_c} \rceil$ blocks.
      \STATE Adjust the allocation of registers according to the determined count of consumer warps.
      \STATE On-chip initialization: set $\mathbf{O}_i = (0) \in \mathbb{R}^{B_r \times d}$, $\ell_i = (0)$, and $m_i = (-\infty) \in \mathbb{R}^{B_r}$.
      \STATE Await transfer of $\mathbf{Q}_i$ and the initial $\mathbf{K}_0$ into shared memory.
      \STATE Using WGMMA, calculate $\mathbf{S}_{\mathrm{cur}} = \mathbf{Q}_i \mathbf{K}_0^T$. Commit and synchronize.
      \STATE Free the buffer stage allocated for the $0$th block of $\mathbf{K}$.
      \STATE Derive $m_{i}$, $\tilde{\mathbf{P}}_{\mathrm{cur}}$, and $\ell_{i}$ using $\mathbf{S}_{\mathrm{cur}}$, and rescale $\mathbf{O}_i$ accordingly.
      \FOR{$1 \le j < T_c - 1$}
        \STATE Await completion of loading $\mathbf{K}_j$ into shared memory.
        \label{code:mainloop_start}
        \STATE Perform WGMMA to compute $\mathbf{S}_{\mathrm{next}} = \mathbf{Q}_i \mathbf{K}_{j}^T$. Commit without waiting. \label{code:first_wgmma}
        \STATE Await transfer of $\mathbf{V}_{j-1}$ into shared memory.
        \STATE Use WGMMA to update $\mathbf{O}_{i} = \mathbf{O}_{i} + \tilde{\mathbf{P}}_{\mathrm{cur}} \mathbf{V}_{j-1}$. Commit and proceed without waiting. \label{code:second_wgmma}
        \STATE Wait for completion of WGMMA $\mathbf{Q}_i \mathbf{K}_{j}^T$.
        \STATE Compute $m_{i}$, $\tilde{\mathbf{P}}_{\mathrm{next}}$, and $\ell_{i}$ based on $\mathbf{S}_{\mathrm{next}}$. \label{code:softmax}
        \STATE After WGMMA $\tilde{\mathbf{P}}_{\mathrm{cur}} \mathbf{V}_{j-1}$ finishes, rescale $\mathbf{O}_i$.
        \STATE Release $(j\,\%\,s)$th and $(j-1\,\%\,s)$th buffer stages corresponding to $\mathbf{K}$ and $\mathbf{V}$, respectively.
        \STATE Assign $\mathbf{S}_{\mathrm{next}}$ to $\mathbf{S}_{\mathrm{cur}}$.
        \label{code:mainloop_end}
      \ENDFOR
      \STATE Wait until $\mathbf{V}_{T_c - 1}$ is accessible in shared memory.
      \STATE Utilize WGMMA to calculate
      $\mathbf{O}_{i} = \mathbf{O}_{i} + \tilde{\mathbf{P}}_{\mathrm{last}} \mathbf{V}_{T_c - 1}$, then commit and synchronize.

\STATE Epilogue: Adjust $\mathbf{O}_{i}$ by scaling according to $m_{i}$. Determine $L_{i}$ utilizing $m_{i}$ and $\ell_{i}$.
Record $\mathbf{O}_{i}$ and $L_{i}$ in HBM, assigning them to the $i$-th segment of $\mathbf{O}$ and $L$.
\end{algorithmic}
\end{algorithm}

\cref{alg:flash3_wgmma} serves as a substitute for the consumer pathway in \cref{alg:flash3_wgmma_ws_only}, thereby completing the \textsc{FlashAttention-3} method for FP16 calculations. Conceptually, WGMMA is used as shorthand to denote asynchronous GEMM operations. During the main loop (lines \ref{code:mainloop_start} through \ref{code:mainloop_end}), the second WGMMA call for iteration $j$ (referenced at line \ref{code:second_wgmma}) is executed concurrently with softmax computations from the subsequent iteration $j+1$ (line \ref{code:softmax}).

Although the pipelined architecture depicted earlier can theoretically boost performance, practical limitations must be accounted for:

\paragraph{Compiler reordering} The presented pseudocode assumes an ideal sequence of operations; however, the NVCC compiler may optimize instruction order through rearrangement. This optimization process can interfere with the intentional alternation of WGMMA and non-WGMMA tasks in the pipeline, thereby reducing or destabilizing performance advantages. Examining the SASS output reveals that, as anticipated, the compiler does produce overlapped instructions (see Section~\ref{sec:2-stage-sass}).

\paragraph{Register pressure} Achieving high efficiency requires avoiding register spills as much as possible. The 2-stage pipeline, however, demands more registers to hold intermediate values and preserve stage information. In particular, one additional $\mathbf{S}_{\mathrm{next}}$ must be retained in registers, incurring extra consumption equivalent to $B_r \times B_c \times \text{sizeof}(\text{float})$ per threadblock. This heightened requirement may constrain choices around larger block sizes, which are themselves intensive in register usage. Profile-guided decisions regarding these compromises are recommended for optimal implementation.

\paragraph{3-stage pipelining} Building on the prior 2-stage method, we introduce a 3-stage approach designed to further overlap the second WGMMA operation with the softmax computation. While this method can increase Tensor Core throughput, it comes at the cost of needing additional registers to support the extra pipeline stage. This increases the complexity of balancing tile size and pipeline depth. Comprehensive information and experimental results for the 3-stage variant are provided in ~\cref{sec:3-stage}.

\subsection{Low-precision with FP8}
\label{sec:algofp8}

\textbf{Efficiency: layout transformations.} Executing the forward pass of \textsc{FlashAttention-3} using FP8 precision introduces specific complications related to layout compatibility that do not arise with FP16.

To begin with, it is important to observe that the input tensors $\mathbf{Q}$, $\mathbf{K}$, and $\mathbf{V}$ are conventionally organized in memory such that the head dimension is contiguous. However, the FP8 WGMMA kernel requires that, for the second GEMM, $\mathbf{V}$ (specifically, the portions of $\mathbf{V}$ transferred into SMEM) be laid out contiguously along the sequence length dimension to comply with the k-major ordering constraint. Because TMA loads cannot alter which dimension is contiguous, the problem must be addressed by either (1) transposing $\mathbf{V}$ within GMEM as a preparatory measure, or (2) performing an in-kernel transpose of segments of $\mathbf{V}$ after these have been loaded into SMEM. For approach (1), one could either (1a) incorporate the transpose into the epilogue of an earlier operation—such as rotary embedding—or (1b) launch a dedicated transpose kernel beforehand, as a pre-processing stage\footnote{A high-performance transpose kernel can operate at speeds close to the device's memory bandwidth~\citep{colfax_cutlass_transpose_2024}.}, thereby swapping the strides of the head and sequence length dimensions. Nonetheless, method (1a) is challenging to accommodate within general-purpose libraries, and approach (1b) incurs excessive overhead in memory-bound scenarios like inference.

For FP8 \textsc{FlashAttention-3}, we instead employ approach (2). In the case of in-kernel transpose, we exploit the capabilities of the LDSM (\verb|ldmatrix|) and STSM (\verb|stmatrix|) instructions, whereby an entire warp of threads cooperatively loads shared memory (SMEM) into register memory (RMEM) and stores from RMEM to SMEM in 128-byte segments.\footnote{As indicated in the PTX documentation, LDSM/STSM carry out transfers of $8 \times 8$ matrices comprising 16-bit elements \cite[\S 9.7.13.4.15-16]{ptx}. For FP8, we pair two 8-bit values into each 16-bit slot, enabling the use of these instructions in this context. However, the transpose variants of LDSM/STSM lack the ability to separate packed 8-bit elements, leading to additional required register operations between LDSM and STSM in order to effectuate the actual tile-level transpose; we leave further specifics aside.} LDSM and STSM are efficient in terms of register usage, which allows their execution within the producer warpgroup, and they provide support for layout transposition during memory transfers. In addition, following the initial iteration, it is possible to arrange the computation such that the transpose of the upcoming $\mathbf{V}$ tile is initiated while the two WGMMAs—which depend on the previous $\mathbf{V}$ tile and the present $\mathbf{K}$ tile—are still running.

Secondly, it is notable that, in contrast to FP16, the arrangement of the FP32 accumulator produced by an FP8 WGMMA diverges from the expected memory layout for operand A stored in registers. Illustrative fragments of these distinct register layouts are provided in \cref{fig:rmem_accum} and \cref{fig:rmem_operand}, where each register entry is presented per thread in sequential order. By employing byte permutation instructions, the accumulator generated by the initial WGMMA can be reorganized into a configuration suitable for the subsequent WGMMA, while also aligning with the register layout of the $\mathbf{V}$ tile resulting from the in-kernel transpose. Concretely, as shown in \cref{fig:rmem_accum}, the sequential register order is transformed to
$$\{ \verb|d0 d1 d4 d5 d2 d3 d6 d7| \},$$
with this register pattern repeated for every block of 8 bytes. Considering the $\mathbf{P}$ tile’s logical structure, such a permutation effectively rearranges its columns (for example, columns $0189$ become the leading four columns). To ensure WGMMA produces the correct result for the output tile, the in-kernel transpose procedure can similarly be designed to emit the $\mathbf{V}$ tile rows in a correspondingly permuted fashion.\footnote{This flexibility introduced by performing the transpose within the kernel obviates the need for shuffle instructions for transferring register content between threads, which was previously required as described in~\citep{colfax_fp8_flashattention_2024}.}

\textbf{Accuracy: block quantization and incoherent processing.} The FP8 (e4m3) format allocates 3 bits for the mantissa and 4 bits for the exponent, yielding a greater degree of numerical error compared to FP16 or BF16. Large-scale models often exhibit extreme outlier values~\citep{dettmers2208llm, sun2024massive} that far exceed the magnitude of the majority of elements, posing challenges for effective quantization. Standard practice involves per-tensor scaling~\citep{micikevicius2022fp8}, where a single scalar is assigned for each tensor (for instance, individually for $\mathbf{Q}$, $\mathbf{K}$, and $\mathbf{V}$).
In order to mitigate numerical inaccuracies in FP8 attention mechanisms, we introduce two specific strategies:

\begin{enumerate}

\item \textbf{Block quantization}: In this method, a single scalar is assigned to represent each block. For tensors $\mathbf{Q}$, $\mathbf{K}$, and $\mathbf{V}$, the input is partitioned into blocks sized $B_r \times d$ or $B_c \times d$, and quantization is then applied to each block individually. This quantization approach can be seamlessly integrated with the operation preceding attention (such as the rotary embedding), without any added latency, given that the rotary embedding is constrained by memory bandwidth rather than computation. Furthermore, since the \textsc{FlashAttention-3} algorithm is inherently block-oriented, each segment of $\mathbf{S}$ can be efficiently rescaled to reflect block quantization without incurring computational overhead.

\item \textbf{Incoherent processing}: To mitigate the effects of outliers, prior to quantization to FP8, we left-multiply both $\mathbf{Q}$ and $\mathbf{K}$ by a random orthogonal matrix $\mathbf{M}$. Since orthogonality implies $\mathbf{M} \mathbf{M}^\top = I$, this guarantees that $(\mathbf{Q} \mathbf{M})(\mathbf{K} \mathbf{M})^\top = \mathbf{Q} \mathbf{K}^\top$, i.e., the attention mechanism’s output remains unchanged. This approach helps diffuse the impact of outliers, as each element in $\mathbf{Q} \mathbf{M}$ or $\mathbf{K} \mathbf{M}$ is constructed as a randomized sum of entries from $\mathbf{Q}$ or $\mathbf{K}$, which acts to lower the quantization error. In line with the procedures of \citet{chee2024quip} and~\citet{tseng2024quip}, we implement $\mathbf{M}$ as a product involving random diagonal matrices with entries $\pm1$ and a Hadamard matrix; this formulation admits multiplication in $O(d \log d)$ operations, as opposed to the $O(d^2)$ required for a general orthogonal transformation. Additionally, this process can be fused with the rotary embedding phase at no increase in computational cost.

\end{enumerate}

Our experiments confirm that these strategies together decrease numerical error by as much as 2.6$\times$, as discussed in \cref{sec:numerical_error}.

\section{Empirical Validation}
\label{sec:exp}

To implement \textsc{FlashAttention-3} and assess its performance and correctness, we leverage CUTLASS~\citep{Thakkar_CUTLASS_2023} primitives, including the WGMMA and TMA abstractions.

\begin{itemize}

\item \textbf{Evaluating attention performance.} We evaluate the execution time of \textsc{FlashAttention-3} over various sequence lengths and benchmark it against several baselines, including the conventional PyTorch attention implementation, \textsc{FlashAttention-2}, \textsc{FlashAttention-2} implemented in Triton (leveraging H100-specific optimizations), and a version of \textsc{FlashAttention-2} from cuDNN, tailored for H100 GPUs. Our experiments demonstrate that \textsc{FlashAttention-3} achieves speedups of up to 2.0$\times$ compared to \textsc{FlashAttention-2} and 1.5$\times$ compared to the Triton variant of \textsc{FlashAttention-2}. Furthermore, \textsc{FlashAttention-3} attains peak throughput of 740 TFLOPs/s, corresponding to 75\% of the H100 GPUs' theoretical peak performance.
\item \textbf{Ablation analysis.} Our studies corroborate that enhancements such as warp-specialization and the integration of GEMM-softmax pipelining are primary drivers for the increased efficiency observed in \textsc{FlashAttention-3}.
\item \textbf{FP8 attention accuracy.} Through empirical assessment, we verify that employing block quantization alongside incoherent computation techniques diminishes the numerical error in FP8 \textsc{FlashAttention-3} by a factor of 2.6$\times$.
\end{itemize}

\subsection{Benchmarking Attention}
\label{subsec:benchmark_attn}

The runtime performance of various attention mechanisms was evaluated on an H100 80GB SXM5 GPU across multiple configurations (both with and without causal masking, and head dimensions set to either 64 or 128), utilizing FP16 input data. The outcomes of these benchmarks are presented in~\cref{fig:benchmark_attn_fwd} and~\cref{fig:benchmark_attn_bwd}. The data illustrate that \textsc{FlashAttention-3} achieves approximately 1.5–2.0$\times$ speedup over \textsc{FlashAttention-2} during the forward computation and is about 1.5–1.75$\times$ faster during the backward computation. When set against a conventional attention implementation, \textsc{FlashAttention-3} demonstrates an acceleration factor ranging from 3$\times$ to 16$\times$. Additionally, for sequence lengths of 1k and higher, \textsc{FlashAttention-3} outperforms a closed-source, vendor-optimized library (cuDNN) designed specifically for H100 GPU architectures.

\paragraph{Benchmark settings:} We vary the sequence length as 512, 1k, ..., 16k, and set batch size so that the total number of tokens is 16k. We set the hidden dimension to 2048, and head dimension to be either 64, 128, or 256 (i.e., 32 heads, 16 heads, or 8 heads). To calculate the FLOPs of the forward pass, we use:

\begin{equation*}
  4 \cdot \text{seqlen}^2 \cdot \text{head dimension} \cdot \text{number of heads}.
\end{equation*}

By applying causal masking, we adjust this value by halving it, reflecting that nearly 50% of the entries are actually computed. For the backward pass FLOPs, we scale the forward pass count by a factor of 2.5; this is because the backward pass involves five matrix multiplications, whereas the forward pass contains only two, a difference stemming from recomputation.

Additionally, we evaluate the forward pass runtime utilizing FP8 within comparable configurations. The outcomes for a headdim of 256 are presented in ~\cref{fig:benchmark_attn_fp8}, and comprehensive results can be found in \cref{sec:benchmark_attn_fp8_full}.

\subsection{Ablation Study: 2-Stage Pipelining Experiments}
\label{sec:2-stage-pipelining-experiments}

We conduct an ablation study by disabling both the two-stage WGMMA-softmax pipeline and warp-specialization in non-causal FP16 \textsc{FlashAttention-3}, holding parameters constant at $\{\text{batch}, \text{seqlen}, \text{nheads}, \text{hdim}\} = \{ 4, 8448, 16, 128\}$.
As shown in~\cref{table:ablation_pipelining}, the results indicate that introducing our proposed techniques—such as asynchronous processing via warp-specialization and overlapping of GEMM with softmax computation—yields notable improvements in throughput, increasing performance from 570 to 661 TFLOPs.

\subsection{Numerical Error Validation}
\label{sec:numerical_error}

Given the growing concern regarding the numerical accuracy~\citep{golden2024flash} of
\textsc{FlashAttention}, we evaluate \textsc{FlashAttention-2}, \textsc{FlashAttention-3}, and a conventional attention mechanism by benchmarking them against a FP64 reference implementation. To mimic the occurrence of outlier features and activations commonly observed in LLMs~\citep{dettmers2208llm,
  sun2024massive}, we sample the elements of $\mathbf{Q}, \mathbf{K}, \mathbf{V}$ according to the distribution below:

\begin{equation*}
  \mathcal{N}(0, 1) + \mathcal{N}(0, 100) \cdot \mathrm{Bernoulli}(0.001).
\end{equation*}

In this setup, each matrix entry follows a normal distribution with a mean of zero and a standard deviation of 1; however, for 0.1\% of the entries, we introduce an additional independent component drawn from a normal distribution with a standard deviation of 10. The root mean squared error (RMSE) outcomes are reported in~\cref{table:numerical_error}. When using FP16, both \textsc{FlashAttention-2} and \textsc{FlashAttention-3} reduce the RMSE by a factor of 1.7 compared to the conventional implementation, primarily because the intermediate softmax values are maintained in FP32 precision. In contrast, the baseline attention implementation in FP8 applies per-tensor scaling, utilizes an FP32 accumulator for the matrix multiplication, and stores intermediate softmax computations in FP16. Owing to block quantization and incoherent computation, \textsc{FlashAttention-3} in FP8 achieves 2.6$\times$ greater numerical accuracy relative to the baseline.

\section{Dicussion, Limitations, Conclusion}
\label{sec:discussion}

Through the development of \textsc{FlashAttention-3}, we have shown that innovative programming methodologies and leveraging hardware capabilities like asynchrony and reduced precision can significantly improve both the speed and fidelity of attention mechanisms. Our approach delivers a 1.5-2.0$\times$ acceleration over \textsc{FlashAttention-2}, and achieves a 2.6$\times$ decrease in FP8 numerical error relative to conventional per-tensor quantization techniques. This study does have certain constraints, including optimizing the design for LLM inference, incorporating a persistent kernel architecture within the FP8 kernel,\footnote{In our experiments, FP16 \textsc{FlashAttention-3} incorporates both persistent kernel and load balancing strategies, whereas FP8 \textsc{FlashAttention-3} currently lacks these features. This is partially why FP8 \textsc{FlashAttention-3} is less effective for short sequence lengths and causal masking, in comparison with FP8 cuDNN kernels.} and further investigating how low-precision attention influences large-scale training dynamics. Although our work primarily utilizes Hopper GPUs, we anticipate that these techniques will be broadly applicable to other accelerator hardware as well. We believe that advancing the efficiency and precision of attention mechanisms can pave the way for novel use cases, particularly within tasks that require long context windows.

\end{document}